\documentclass[11pt,letterpaper]{letter}
\usepackage[margin=1in]{geometry}
\usepackage{hyperref}

\signature{Michael Whiteside \\ Independent Researcher \\ United Kingdom}
\address{Michael Whiteside \\ Independent Researcher \\ United Kingdom \\ mickwh@msn.com}

\begin{document}

\begin{letter}{Editorial Office \\ PLOS Computational Biology}

\opening{Dear Editors,}

I am pleased to submit the manuscript ``Temporal Persistence as a Measurable Property of Gene Expression: AR(2) Eigenvalue Analysis Reveals Circadian Hierarchy, Phase-Gated Regulation, and Fibonacci-Like Dynamics Across Mammalian Tissues'' for consideration as a Research Article in \textit{PLOS Computational Biology}.

This manuscript presents a unified computational framework---grounded in second-order autoregressive (AR(2)) modelling---for quantifying temporal persistence in gene expression. Starting from the biological observation that colonic stem cell renewal involves generational memory, we develop and validate an analytical pipeline that, from just two fitted coefficients, recovers the known circadian gene hierarchy across 12 mouse tissues and 4 species (mouse, human, baboon, \textit{Arabidopsis}).

The paper makes five principal contributions:

\begin{enumerate}
\item The eigenvalue modulus $|\lambda|$ from AR(2) regression blindly recovers clock $>$ target $>$ background gene hierarchy (clock median $|\lambda| = 0.647$, $p < 0.001$), validated across 20,955 genes, 12 tissues, and 4 species.

\item Phase-gated extension (PAR(2)) identifies Wee1 as gated by all 8 clock genes across 4--6 tissues each---the strongest circadian hub in the cell cycle network---with cross-tissue FDR $\sim$2\%.

\item Systems-level eigenperiod analysis reveals a robust healthy--cancer separation: 7--13h ultradian dynamics in healthy tissues versus 22--23h near-circadian periods in cancer models ($p < 10^{-15}$).

\item In tissues requiring precise temporal coordination, AR(2) coefficient ratios cluster near the golden ratio $\phi \approx 1.618$ with 48-fold enrichment, connecting temporal dynamics to Fibonacci-like recursion.

\item The framework resolves two long-standing paradoxes in colonic crypt renewal and generates four falsifiable predictions with quantitative success criteria.
\end{enumerate}

The manuscript includes comprehensive validation: 12-analysis robustness suite, ODE model bridge testing, three automated bias audits, BMAL1-knockout causal validation, cross-metric independence testing, and clinical trial cross-validation (PALOMA-3, 302 patients).

All analyses are fully reproducible via our open-source platform (\url{https://par2-discovery-engine.replit.app}) and GitHub repository (\url{https://github.com/mickwh2764/par2-discovery-engine}).

This work is original, has not been published elsewhere, and is not under consideration at another journal. A companion manuscript on Fibonacci-like dynamics in circadian gene expression is under review at \textit{The Fibonacci Quarterly}.

I believe this manuscript is well-suited for \textit{PLOS Computational Biology} because it introduces a validated computational methodology with broad biological applications, supported by extensive cross-species and cross-tissue validation.

\closing{Sincerely,}

\end{letter}
\end{document}
